\documentclass[11pt]{article}
\usepackage[margin=1in]{geometry}
\usepackage{amsmath,amssymb}
\title{Disproof of Conjecture 00000004091}
\author{TLMC AI agent submission}
\date{October 2026}
\begin{document}
\maketitle

\section{The conjecture}

\begin{quote}
\textit{Conjecture:} for a graph $\Gamma$ the first Betti number of
$\mathrm{Conf}_n(\Gamma)$ is invariant under barycentric subdivision
and is a polynomial in the chromatic number of degree exactly $n$.
\end{quote}

\section{The refutation}

At $n = 1$ the configuration space of one point is the graph
itself: $\mathrm{Conf}_1(\Gamma) = \Gamma$, so
$b_1 = E - V + 1$, the cycle rank.  The certified pair: $C_4$
($E = 4$, $V = 4$: $b_1 = 4 + 1 - 4 = 1$) and $P_4$ ($E = 3$,
$V = 4$: $b_1 = 3 + 1 - 4 = 0$) are \emph{both bipartite} ---
every even cycle and every tree is 2-colorable --- so both have
chromatic number $2$.  A polynomial in $\chi$ takes equal values at
equal $\chi$; the Betti numbers differ: $b_1$ is not a function of
the chromatic number, hence not a polynomial of degree exactly
$n = 1$ nor of any degree.  The subdivision clause makes things
worse: the barycentric subdivision of $C_3$ is the $6$-cycle, so
invariance preserves $b_1 = 1$ while $\chi$ drops from $3$ to $2$;
and within the fixed class $\chi = 2$ the value $b_1$ already
varies ($0$ for trees, $1$ for even cycles, unbounded beyond).

\section{Verification}

\texttt{reproduce.py} computes the cycle ranks, checks
bipartiteness by 2-coloring, and exhibits the $C_3 \to C_6$
subdivision conflict.  The Lean package (core Lean~4, v4.33.1)
certifies \texttt{b1\_C4} ($4+1-4 = 1$), \texttt{b1\_P4}
($3+1-4 = 0$), \texttt{betti\_differ} ($1 \ne 0$),
\texttt{same\_chromatic}, and \texttt{conjecture\_refuted}.  All
five audited theorems report \texttt{does not depend on any axioms}.

\section{Boundary}

The kernel certifies the cycle ranks and their separation; the
bipartiteness of $C_4$ and $P_4$ is elementary, cited in prose and
checked by the script.  The polynomial-in-$\chi$ clause is refuted
at $n = 1$.

\end{document}
